\begin{lemma}[Attainment of the Tashkinov order]
Let $G$ be a finite loopless multigraph satisfying the tablet's
edge-criticality condition.  If $\chi'(G)=k+1$ and
$\Delta(G)+1\leq k$, then there are an edge identity $e_0$, a proper
partial $k$-edge-coloring $\varphi$ leaving exactly $e_0$ uncolored, and a
corresponding Tashkinov tree $T$ such that
$|V(T)|=t(G)$.  Equivalently, $(T,e_0,\varphi)$ is globally maximum.
\end{lemma}
\begin{proof}
First, $E(G)$ is nonempty.  If it were empty, the unique map from its edge
type to partial colors in $\operatorname{Fin}(0)$ would be both proper and
full, vacuously satisfying \noderef{IsProperEdgeColoring} and
\noderef{HasKEdgeColoring}.  The minimum characterization
\noderef{ChromaticIndexSpec} would then give $\chi'(G)=0$.  On the other
hand, $\Delta(G)+1\leq k$ gives $1\leq k$, and the assumed equality gives
$\chi'(G)=k+1\geq2$, a contradiction.  Choose an edge identity $e_0$.

By \noderef{Critical}, deleting $e_0$ strictly lowers the chromatic index,
so
\[
 \chi'(G-e_0)<\chi'(G)=k+1,
 \qquad\text{hence}\qquad \chi'(G-e_0)\leq k.
\]
The existence half of \noderef{ChromaticIndexSpec} supplies a proper full
edge-coloring of $G-e_0$ with colors in
$\operatorname{Fin}(\chi'(G-e_0))$.  Embed this finite color set into
$\operatorname{Fin}(k)$ using the displayed inequality.  On every retained
edge use the embedded color, and give $e_0$ the value ``uncolored.''
Call the resulting partial coloring $\varphi$.  Two distinct colored
adjacent edge identities are retained identities of $G-e_0$, so their
preimages had different colors and their embedded colors are still
different.  Thus $\varphi$ is proper.  It is uncolored precisely on $e_0$,
as required by \noderef{IsOneEdgeUncolored}.  Together with the three
numerical hypotheses, this is exactly the $k$-triple condition
\noderef{KTriple}.

The equality $\chi'(G)=k+1$ also gives $\chi'(G)-1=k$.  Thus the color
parameter in the attainable-order set defining \noderef{TashkinovOrder}
is exactly the $k$ used by the triple just constructed.

Write the unordered endpoint pair of $e_0$ as $\{x,y\}$.  The loopless law
in \noderef{LooplessMultigraph} gives $x\ne y$.  Form the length-one tree
sequence with ordered vertices $x,y$ and sole edge identity $e_0$.  The
vertices are distinct, the sole edge attaches $y$ to the earlier vertex
$x$, and the first-edge clause is immediate, so this is a
\noderef{TreeSequence}.  Its \noderef{TashkinovTree} condition is vacuous,
because there is no tree edge after the first one.  Consequently the set
of attainable orders used in \noderef{TashkinovOrder} is nonempty; it
contains the order of this two-vertex initial tree.

For completeness, maximum attainment is not being inferred merely from
the word ``maximum.''  Every tree sequence has pairwise distinct vertices
by \noderef{TreeSequence}, so \noderef{TreeSequenceVertexSet} has
cardinality at most $|V(G)|$.  Hence all attainable orders lie in the
finite set $\{0,1,\ldots,|V(G)|\}$, and the preceding initial tree makes
the attainable subset nonempty.  The finite set therefore has a largest
member $m$, and \noderef{TashkinovOrder} defines $t(G)$ to be that member.
Membership of $m$ in the attainable-order set yields witnesses
$e$, $\psi$, and $S$ with the $k$-triple condition, the Tashkinov-tree
condition, and $|V(S)|=m=t(G)$.  These three clauses are precisely
\noderef{GlobalMaximumTriple}, so $(S,e,\psi)$ is the required globally
maximum triple.
\end{proof}
